% SEGMENT OLP-0412-S01
% Part: normal-modal-logic
% Chapter: syntax-and-semantics
% Section: relational-models

\documentclass[../../../include/open-logic-section]{subfiles}

\begin{document}

\olfileid{nml}{syn}{rel}

\olsection{தொடர்புசார் மாதிரிகள்}

இயல்பான வாய்ப்புநிலைத் தருக்கங்களின் பொருண்மையியலுக்குரிய அடிப்படைக்
கருத்து ஒரு \emph{தொடர்புசார் மாதிரி} ஆகும். இருநிலை ``அணுகுறவால்''
தொடர்புபடுத்தப்பட்ட உலகங்களின் கணமும், எந்தெந்த உலகங்களில் எந்தெந்த
!!{propositional variable}s ``மெய்'' எனக் கொள்ளப்படுகின்றன என்பதைத்
தீர்மானிக்கும் ஒரு மதிப்பளிப்பும் இதில் அடங்கும்.

\begin{defn}
  அடிப்படை வாய்ப்புநிலை மொழிக்கான ஒரு \emph{மாதிரி} என்பது $\mModel{M}
  = \tuple{W, R, V}$ என்ற மும்மை; இங்கு
  \begin{enumerate}
  \item $W$ என்பது ``உலகங்களைக்'' கொண்ட ஒரு வெறுமையற்ற கணம்;
  \item $R$ என்பது $W$ மீதான ஓர் இருநிலை அணுகுறவு; மேலும்
  \item $V$ என்பது ஒவ்வொரு !!{propositional variable}~$p$ க்கும்
    சாத்தியமான உலகங்களின் கணம் $V(p)$ ஐ மதிப்பளிக்கும் ஒரு சார்பு.
  \end{enumerate}
  $Rww'$ மெய்யாகும்போது, $w$ இலிருந்து $w'$ ஐ \emph{அணுகலாம்} என்கிறோம்.
  $w \in V(p)$ ஆகும்போது, $w$ இல் $p$ \emph{மெய்} என்கிறோம்.
\end{defn}

எளிய வரைபடங்களைக் கொண்டு மாதிரிகளைச் சித்தரிக்க முடிவதே தொடர்புசார்
பொருண்மையியலின் பெரும் நன்மை; \olref{fig:simple} இல் உள்ள வரைபடம் இதற்கு
ஓர் எடுத்துக்காட்டு. உலகங்கள் முனைகளால் குறிக்கப்படுகின்றன; $w$ இலிருந்து
$w'$ க்கு ஓர் அம்பு இருந்தால், எனில் மட்டுமே, உலகம்~$w$ இலிருந்து
உலகம்~$w'$ ஐ அணுகலாம். மேலும், $w \in V(p)$ ஆகும்போது ஒரு முனைக்கு
(உலகத்துக்கு)~$\mTrue{p}$ என்றும், இல்லையெனில்~\mFalse{p} என்றும்
குறியிடுகிறோம். $W = \{w_1, w_2, w_3\}$, $R = \{\tuple{w_1,
  w_2},\tuple{w_1, w_3}\}$, $V(p) = \{w_1, w_2\}$, $V(q) = \{w_2\}$
ஆகியவற்றைக் கொண்ட மாதிரியை \olref{fig:simple} சித்தரிக்கிறது.

\begin{figure}
  \begin{center}
    \begin{tikzpicture}[modal]
      \node[world] (w1) [label={[align=right]right:\mTrue{p}\\ \mFalse{q}}]{$w_1$};
      \node[world] (w2) [label={[align=right]right:\mTrue{p}\\ \mTrue{q}},
        above right=of w1]{$w_2$};
      \node[world] (w3) [label={[align=right]right:\mFalse{p}\\ \mFalse{q}},
        below right=of w1] {$w_3$};
      \draw[->] (w1) to (w2);
      \draw[->] (w1) to (w3);
    \end{tikzpicture}
  \end{center}
  \caption{ஓர் எளிய மாதிரி.}
  \ollabel{fig:simple}
\end{figure}

\end{document}
